%----------
%	HATE SPEECH DETECTION
%----------	

In Sections~\ref{sec:hs_society},~\ref{sec:hs_sports} and~\ref{sec:hs_socialmedia} we have seen that hate speech is a common problem in our everyday lives, especially on social media. Historically, the identification of hate speech was performed in a manual way, where human moderators would review and detect offensive content. However, with the exponential growth of social media, manual detection has proven to be insufficient, as it is time-consuming and often inconsistent and biased \cite{sap-2019}. To address these challenges, the focus has increasingly shifted towards automated solutions in the Artificial Intelligence landscape, as illustrated in Figure~\ref{fig:ai_landscape}. In particular NLP solutions using Machine Learning, Deep Learning and, more recently, Generative AI.

\begin{figure}[H]
	\ffigbox[\FBwidth]
	{\caption{Diagram of the AI landscape, adapted from \url{https://neosnetworks.com/resources/blog/what-is-ai-in-networking/}}\label{fig:ai_landscape}} 
 % [Aquí ponemos como queremos que aparezca en el índice, sin la referencia]{Aquí como queremos que aparezca en el documento, con la referencia}
	{\includegraphics[scale=0.5]{Plantilla_TFG_ingles_2019/imagenes/TFG Pipeline-AI-landscape.drawio.pdf}}
\end{figure}

One crucial challenge is that the majority of the studies that have been done in this field are in English, due to the availability of vast amounts of resources in this language. Thus, studies that attempt to research in other languages face many problems during their investigations. In particular, \cite{fersini-2018} presents resources for the identification and classification of misogynistic language on Twitter for Italian and Spanish. They also were the organizers of IberEval - Automatic Misogyny Identification (AMI) competition\footnote{From \url{https://sites.google.com/view/ibereval-2018/home?authuser=0}, last visited 17/08/24}.

In Spanish, the SINAI group of the Universidad of Jaen is a research group that has conducted different studies in the field of natural language processing and is a reference in the field of detecting HS in Spanish. In particular, they have created accurate language resources in Spanish for text-based HS detection of xenophobia and misogyny in tweets \cite{plazaDelArco-2020}, while in \cite{plazaDelArco-2021} the focus was on comparing existing pre-trained language models in Spanish.

Additionally, as mentioned in Section~\ref{sec:hs_socialmedia}, Hatemedia developed the "Monitor de Odio", which provides a comprehensive indicator of the presence, intensity, and types of hate expressions in Spanish media. This tool was trained using a dataset of 1,100,742 messages collected from Spanish digital media in January 2021, with 324,395 unique messages being analyzed. It specifically focuses on content from five major Spanish newspapers: 20 Minutos, ABC, El Mundo, La Vanguardia, and El País. The Monitor de Odio classifies hate speech based on its presence, type (general, political, sexual, misogynistic, or xenophobic), and intensity (levels 1 to 4), as described in Section~\ref{sec:hs_socialmedia}. To develop these three tasks, a variety of data preprocessing techniques were employed, including the removal of null values and duplicates, elimination of URLs, emojis, and mentions, removal of empty rows, conversion of all text to lowercase, stopword removal, and tokenization. These processed data were then used to develop and apply multiple machine learning and deep learning models across three distinct tasks.

The first one is \cite{degregorio-2024}, “Algorithm for classifying hate expressions in Spanish”\footnote{From \url{https://github.com/esaidh266/Algorithm-for-detection-of-hate-speech-in-Spanish}, last visited 17/08/24}. This work aimed to classify content into two categories: "Odio" (eng. `hate') and "No Odio" (eng. `not hate'), and proposes three different train-test splits (90-10, 80-20, and 70-30) for the evaluation of the performance of various machine learning models, including naive Bayes, Random Forest, Gradient Boosting, MLP, CART, and SVM. The best results were achieved using SVM, Gradient Boosting, MLP, and CART, with CART being selected as the top performer, achieving an accuracy of over 97\%. 

The second task is \cite{blanco-2024}, “Algorithm for classifying hate expressions by intensities in Spanish”\footnote{From \url{https://github.com/esaidh266/Algorithm-for-classifying-hate-expressions-by-intensities-in-Spanish}, last visited 17/08/24}. This study aimed to classify content by hate intensity on a scale from 1 to 6 (which was later reduced to a scale of 1-4 during the methodology phase). The models used in this work include naive Bayes, Gradient Boosting, Random Forest, MLP, and SVM. In addition to these traditional ML models, more complex deep learning models such as RoBERTuito \cite{perez-2023} and LSTM were also implemented. After extensive testing, the CART model, combined with the ADASYN balancing technique and an 80-20 training split, achieved an accuracy of 71.2\% and was selected as the best-performing model.

The third one is \cite{perez-2024}, “Algorithm for classifying hate expressions by types in Spanish”\footnote{From \url{https://github.com/esaidh266/Algorithm-for-classifying-hate-expressions-by-type-in-Spanish}, last visited 17/08/24}. This approach evaluated the same ML models as in the previous tasks but for classifying the type of hate speech (general, political, sexual, misogynistic, or xenophobic). Along with RoBERTuito, they used different data balancing techniques (SMOTE, ADASYN, Random Oversampling) and train-test splits (90-10, 80-20, and 70-30). Finally, the RoBERTuito model achieved the highest accuracy (76.9\%), making it the best approach among those tested.

Other important sources of HS detection development and research are academic events and shared tasks, where researchers collaborate, share findings, and benchmark their approaches. Notable among these is the AMI (Automatic Misogyny Identification) task\footnote{From \url{https://sites.google.com/view/ibereval-2018/}, last visited 19/08/24}, which focuses on detecting misogynistic content on Twitter in both Spanish and English. It consists of two subtasks: "Task A: Misogyny Identification", which is a binary classification as either misogynous or not misogynous; and "Task B: Misogynistic Behaviour and Target Classification", which target is to classify the misogynous tweets according to both the misogynistic behaviour ('Stereotype \& Objectification', 'Dominance', 'Derailing', 'Sexual Harassment \& Threats of Violence' and 'Discredit') and the target of the message ('Active' and 'Passive').

Similarly, the EXIST (sEXism Identification in Social neTworks)\footnote{From \url{http://nlp.uned.es/exist2024/}, last visited 19/08/24} series aims to capture a broad spectrum of sexist behaviours on social media, defining six different tasks. It includes "Task 1: Sexism Identification in Tweets", which involves a binary classification of tweets as either 'Sexist' or 'Not Sexist'. For tweets identified as 'Sexist', "Task 2: Source Intention in Tweets" classifies them based on the author's intent in 'Direct', 'Reported', or 'Judgemental'. "Task 3: Sexism Categorization in Tweets" further analyzes these tweets to identify which aspects of women are most frequently targeted, in 'Ideological \& Inequality', 'Stereotyping \& Dominance', 'Objectification', 'Sexual Violence', and 'Misogyny \& Non-Sexual Violence'. Moreover, this analysis is extended to visual content with "Task 4: Sexism Identification in Memes", which determines whether a meme is sexist, followed by "Task 5: Source Intention in Memes", which categorizes memes based on the creator's intent, and "Task 6: Sexism Categorization in Memes", which applies the same categories used in Task 3 to classify sexist memes.

Additionally, the Workshop on Online Harms and Abuse (WOAH)\footnote{From \url{https://www.workshopononlineabuse.com/}, last visited 19/08/24}, provides a forum for interdisciplinary research to address topics such as developing new detection models, understanding the biases and limitations of current systems, creating new datasets, and examining the social, legal, and ethical implications of content moderation and online abuse. Now in its eighth edition and co-located with NAACL 2024 in Mexico City, this year's workshop target is on the challenges posed by LLMs in generating toxic content \cite{woah-2024}.

Finally, Generative AI Models have recently gained popularity, especially following the release of ChatGPT\footnote{From \url{https://openai.com/index/chatgpt/}, last visited 07/08/24} in November 2022. In particular, GPT-3.5 is a large language model developed by OpenAI and trained on extensive amounts of data to process and generate text in different domains and for many tasks, one of which could be detecting hate speech \cite{zhang-2023}. \cite{pendzel-2023} examined the performance of various complex deep learning models such as BERT, RoBERTa, and ALBERT, against fine-tuned LLMs like GPT-3.5 for zero-shot hate detection. The results revealed that GPT-3.5, when evaluated on their test sets, "achieved the highest recall and F1 score across all the methods tested". Moreover, \cite{plazaDelArco-2021} assessed the performance of recent pre-trained language models based on the Transformer mechanism, using both multilingual models (BERT, XLM) and a monolingual model (BETO) designed specifically for Spanish. Their findings indicated that the monolingual BETO model outperformed the multilingual models mBERT and XLM. These research directions suggest that pre-trained language models and generative models are valuable tools in identifying hate speech and abusive language. 


